\documentclass{article}
\usepackage[spanish]{babel}
\usepackage[backend=biber,style=numeric]{biblatex} % Cambia 'apa' por 'numeric'
\addbibresource{references.bib} % Indica el nombre de tu archivo .bib
\usepackage{hyperref} % Permite hipervínculos

\begin{document}
\title{Aporte de Turing a la Derrota del Nazismo}
\author{Carla Sumani Pérez Valera C-512  \\ Francisco Vicente Suárez Bellón C-512}
\date{2024}

\maketitle
\section*{Resumen}
En el contexto de uno de los episodios más significativos
de la historia humana, la Segunda Guerra Mundial, se evidenció
la crucial función de las máquinas en la resolución de problemas
complejos que resultan inalcanzables para el ser humano.
El desciframiento de los mensajes encriptados por los \textbf{Nazis}
a través de la \textbf{Máquina Enigma} fue fundamental 
para poner fin al conflicto. Los esfuerzos conjuntos de 
naciones de los Aliados como Polonia y el Reino Unido para 
formar un equipo dedicado a desentrañar estos mensajes, 
junto con la figura destacada de \textbf{Alan Turing } en
el diseño y construcción de la máquina que logró desencriptar
a Enigma, jugaron un papel esencial en el desarrollo
de los acontecimientos.


\section*{Palabras Clave}
\begin{itemize}
    \item Máquina Enigma
    \item Código Enigma
    \item  Nazismo
    \item Criptografía
\end{itemize}

\section*{Contexto Histórico}

Para comenzar a hablar sobre el aporte de Alan Turing en el desciframiento 
del código Enigma y sus repercusiones futuras en las Ciencias de la Computación 
debemos situarnos en el contexto Histórico donde se desarrollaron los hechos.
\linebreak % Espacio entre la aclaracion de pq se hace el resumen histórico y el resto
Con la Invación a Polonia el 1ero de septiembre de 1939 se 
dio inicio a la Segunda Guerra Mundial
por parte de la Alemania dirigida por \textbf{Adolf Hitler} hasta la 
rendicion del Imperio Japones el 2 de septiembre de 1945\cite{Wikipedia:WorldWarIIWiki}.
En ella existían dos principales bloques beligerantes: Fascistas y Aliados, donde sus mayores representantes
fueron, Alemania, Italia y Japón por parte del eje Fascista y La Union Soviética, Estados Unidos de América, China, Francia, Reino Unido y Polonia.
Este fue el conflicto bélico el cual tiene el título por ser el mayor enfrentamiento 
del siglo XX y el más desbastador de la historia humana. En el cual se desarrollaron 
algunas de las batallas mas relevantes de la historia humana como: Kursk, Stalingrado, 
Leningrado, Okinawa, Invación de Normandía o la 
Toma de Berín,
Dunkerque, además de algunos de los mayores horrores de la humanidad como 
el Holocausto o el Bombaedeo de armas Atómicas a las ciudades japonesas de 
Hiroshima y Nagasaki, así como la muerte de mas de 60 millones de personas \cite{Wikipedia:WorldWarIIWiki}. Dado la intensidad del conflicto así como
la fortaleza de los países involucrados la competencia entre ambos bandos
por el desarrollo tecnológico como el Vergeltungswaffe 2 % CItar v2
El sistema de Radar y la Máquina Enigma, esta última
vital para las comunicaciones secretas del ejercito del Führer

Dado lo vital que era para el bando Aliado conocer los movimientos 
de Berlín fue necesario poder decodificar el sistema empleado por
los Alemanes para el envio de mensajes.


\subsection{La historia de la Máquina Enigma}

Dado que ya desde la Primera Guerra Mundial, Reino Unido 
habia interceptado mensajes secretos del bando teuton, como en 
1916 cuando el comandante de la operación de señales alemana en Oriente 
Medio envió un mensaje borracho a todos sus operadores deseándoles una 
Feliz Navidad. Con poca otra actividad ocurriendo durante el 
período navideño, el mismo mensaje aislado y claramente 
idéntico se envió en seis códigos diferentes, solo uno de 
los cuales, hasta ese momento, los británicos habían logrado 
romper. % Añadir referencia de colossus the secrets....
Se evidencia los esfuerzos de ambos bandos 
tanto en las tecnicas de cifrado, por los alemanes,
como los intentos de desencriptar de los britanicos. 



La máquina Enigma fue patentada por el inventor alemán 
\textbf{Arthur Scherbius} al final de la Primera Guerra Mundial, 
concretamente en 
 1918. Scherbius buscaba modernizar la seguridad en las 
 comunicaciones, y su dispositivo electromecánico reemplazó 
 los sistemas de lápiz y papel usados hasta entonces. La idea 
 de Scherbius era crear una máquina que pudiera cifrar y 
 descifrar mensajes de forma rápida y segura, utilizando 
 tecnología de vanguardia para la época. % En el colossus the secrets....

Su uso militar no fue hasta 1926 donde se adopto por la Armada 
de la República de Weimar, la cual capaz de cifrar y descifrar 
mensajes fue muy util en esta campo por lo cual 2 años más tarde en 1928
fue utilizada por el Ejercito de dicho país. Su robustez criptografica seguido por su diseño compacto y 
movil acompañado de su facil empleo, fue ideal para los Nazis para su empleo en la Guerra
para enviar mensajes de alta sensibilidad como avisar de ataques, posicionamientos
y otros detalles de la contienda bélica. La máquina necesitaba de 
una configuracion para encriptar y desencriptar mensajes, donde ambas maquinas
la de transmicion y recepcion de los datos tenian que tener iguales configuraciones.
Un aspecto a tener en cuenta de esta máquina es su capacidad para codificar 
la información en alfabetos extremadamente grandes, haciendo la intercepción indebida
hasta la fecha realizada manualmente por personas, imposible dado que no era factible
lograr las traducciones al alfabeto normal en un tiempo razonablemente humano; por los factores anteriores
el ejercito liderado por Hitler confió en este sistema. 
Los mensajes una vez  Traducidos  serián escritos en codigo Morse para ser 
transmitodo de forma inalámbrica, lo cual no traía vulnerabilidades a los alemanes por su seguridad.

Existián en varios sectores, asi como distintas versiones de esta como el de la  Wehrmacht los utilizados en la Marina
por los U-Boots, sienda esta en la quie se enfocaría Turing por descrifrar. Algunos de los componentes de ellas debian ser reconfigurados cada ciertos periodos de tiempo, que variaban dependiendo
del rotor, fuerza militar y momento de la guerra aumentanto la frecuencia de estos, con el fin de aumentar 
la complejidad a los adversarios para su rotura, esto le dio una clara ventaja al inicio de la Guerra
para desarrollar sus principios de guerra rapida (blitzkrieg). % Anadir ref wiki 



\subsubsection{Hablando sobre la máquina}

  

El artilugio tiene una apariencia muy similar con una máquina de escribir de la época,
en ella se diferencian
los siguientes componentes de interés:
\begin{itemize}
    \item El teclado de escritura
    \item  Rotores (o Ruedas)
    \item Tablero de Conexiones (Steckerbrett)
    \item  Reflector (Umkehrwalze)
    \item Ajustes de la Máquina (Ringstellung)
    \item Teclado de luces (lampboard)
    \item  Indicador
\end{itemize}
 
% Insertar foto de la enigma

En su funcionamiento al ser presionada una tecla en el teclado de escritura un
impulso electrico viaja desde la 
tecla presionada hasta su análoga en la tabla de conexiones, de donde se redirecciona
esta a los rotores, inicialmente configurados, donde por su propia configuración
una letra x podia transformarse en una letra y, además que pulsar repetidamente la misma letra
hacia que se trajera en otra distinta de la anterior finalmente en el teclado de luces encriptada.
Haciendo que teniendo la máquina debidamente configurada escribir el mensaje encriptado anteriormente
en esta hacia salir el mensaje real por el teclado de luces.

Como se ha comentado antes la unica forma de desencriptar un mensaje es de partiendo de la configuracion
inicial de la máquina con la que fue encriptada, en dicha configuración esta involucrado el Tablero de Conexiones
y Los Rotores


\subsubsection{Tabla de conexiones}
El tablero de conexiones (Steckerbrett) ,
era un componente fundamental de la máquina Enigma que añadía una capa adicional
de complejidad al proceso de cifrado. Su función principal era intercambiar pares
de letras antes de que la corriente eléctrica pasara a los rotores, lo que aumentaba
significativamente la seguridad del cifrado.

El Steckerbrett estaba situado en la parte frontal de la máquina y consistía
en un panel con 26 sockets, una para cada letra del alfabeto. 
El operador utilizaba cables con conectores dobles para unir pares de letras en el tablero.
Ejemplo: Si el operador conectaba las letras  A  y  C  en el tablero de conexiones,
 al presionar la tecla  A  en el teclado, la corriente eléctrica primero se dirigía
  a la posición  C  en el Steckerbrett. A partir de ahí, la corriente seguía su camino
   normal a través de los rotores y el reflector, como si se hubiera presionado la tecla  C .
El tablero de conexiones aumentaba la seguridad de Enigma de varias maneras:
\begin{itemize}
    \item Mayor número de combinaciones: Cada cable en el Steckerbrett
    creaba una permutación adicional en el proceso de cifrado.
    Con diez cables conectados, el número de posibles configuraciones del
    Steckerbrett era astronómicamente grande.
    \item Ofuscación de las relaciones entre letras: 
    El intercambio de letras en el Steckerbrett dificultaba 
    el análisis de frecuencia de letras, una técnica comúnmente 
    utilizada por los criptoanalistas para descifrar códigos.
    \item Carácter recíproco: La conexión de pares de letras en el Steckerbrett
    era recíproca. Esto significa que si A se conectaba a C, entonces C también se conectaba a  A . 
    Aunque esta característica simplificaba la operación para los operadores alemanes, también representaba 
    una debilidad que los criptoanalistas aliados pudieron explotar.
\end{itemize}

\subsubsection{Rotores}

Los rotores, también conocidos como ruedas, eran el corazón del mecanismo de cifrado de la máquina Enigma. 
Su función principal era realizar una serie de sustituciones alfabéticas complejas que convertían la letra 
presionada en el teclado en una letra diferente, que luego se iluminaba en el tablero de lámparas.

Aquí se describe el funcionamiento de los rotores en detalle:

\begin{enumerate}
    \item \textbf{Estructura del Rotor:}
    \begin{itemize}
        \item \textbf{Núcleo cableado:} Cada rotor consistía en un disco con 26 contactos eléctricos 
        en cada lado, uno para cada letra del alfabeto. En el interior del disco, los contactos de un 
        lado estaban conectados a los del otro lado mediante un intrincado cableado. Este cableado era 
        único para cada rotor y determinaba la sustitución alfabética que realizaba el rotor.
        \item \textbf{Anillo exterior y muesca:} El rotor tenía un anillo exterior que se podía girar en 
        relación con el núcleo cableado. Este anillo tenía una muesca (o varias muescas en algunos rotores) 
        que desempeñaba un papel crucial en el mecanismo de rotación de los rotores.
    \end{itemize}
    
    \item \textbf{Movimiento Rotatorio:}
    \begin{itemize}
        \item \textbf{Avance del rotor derecho:} Cada vez que se presionaba una tecla en el teclado, 
        el rotor de la derecha avanzaba una posición. Este movimiento rotatorio era similar al de un 
        odómetro.
        \item \textbf{Turnover:} Cuando la muesca del anillo exterior del rotor derecho se alineaba 
        con un punto fijo en la máquina, el rotor del medio avanzaba una posición. Este mecanismo, 
        conocido como "turnover", introducía una mayor complejidad al cifrado, ya que la sustitución 
        alfabética no solo cambiaba con cada pulsación de tecla, sino que también cambiaba a medida 
        que los rotores giraban.
        \item \textbf{Múltiples turnovers:} En máquinas Enigma más avanzadas, algunos rotores tenían 
        dos muescas, lo que podía provocar que el rotor de la izquierda también avanzara en ciertos puntos, creando un ciclo de rotación aún más complejo.
    \end{itemize}

    \item \textbf{Interacción con otros componentes:}
    \begin{itemize}
        \item \textbf{Tablero de conexiones (Steckerbrett):} Antes de llegar a los rotores, la corriente eléctrica pasaba por el tablero de conexiones, donde se intercambiaban pares de letras. Este proceso de intercambio añadía otra capa de sustitución al cifrado.
        \item \textbf{Reflector (Umkehrwalze):} Después de pasar por los tres rotores, la corriente eléctrica llegaba al reflector, que la enviaba de vuelta a través de los rotores por una ruta diferente. Este proceso aseguraba que una letra nunca se cifrara como sí misma.
    \end{itemize}

    \item \textbf{Importancia para la seguridad:}
    \begin{itemize}
        \item \textbf{Combinaciones casi ilimitadas:} La combinación de diferentes rotores, sus posiciones iniciales, la configuración del Steckerbrett y el mecanismo de turnover generaba un número astronómicamente grande de combinaciones de cifrado posibles.
        \item \textbf{Sustituciones dinámicas:} El movimiento rotatorio de los rotores creaba una secuencia de sustituciones alfabéticas que cambiaba constantemente, lo que dificultaba enormemente el descifrado del código mediante análisis de frecuencia u otras técnicas convencionales.
    \end{itemize}
\end{enumerate}

%  "Copeland Enigma.pdf".

\subsection{Modificaciones}


Enigma experimentó una evolución significativa a lo largo del tiempo 
para aumentar su seguridad y dificultar la labor de los criptoanalistas aliados. 
A continuación, se describen algunos aspectos clave de su evolución:

\begin{enumerate}
    \item \textbf{Complejidad Creciente del Tablero de Conexiones (Steckerbrett):}
    \begin{itemize}
        \item Al principio de la guerra, el Steckerbrett solía conectar solo seis pares de letras.
        \item A medida que avanzaba la guerra y los criptoanalistas aliados lograban descifrar 
        el código Enigma con mayor frecuencia, los alemanes aumentaron el número de letras conectadas 
        en el Steckerbrett para incrementar la complejidad del cifrado. Este cambio hizo que el descifrado 
        fuera mucho más difícil, ya que aumentaba exponencialmente el número de posibles configuraciones del 
        Steckerbrett.
    \end{itemize}

    \item \textbf{Introducción de Rotores Adicionales:}
    \begin{itemize}
        \item Inicialmente, la máquina Enigma utilizaba tres rotores para realizar las sustituciones 
        alfabéticas.
        \item Posteriormente, se introdujeron máquinas Enigma con cuatro rotores, lo que añadía una capa 
        adicional de complejidad al cifrado. La inclusión de un cuarto rotor aumentaba significativamente 
        el número de posibles configuraciones, haciendo que la tarea de descifrado fuera mucho más difícil.
    \end{itemize}

    \item \textbf{Cambios en los Patrones de las Ruedas:}
    \begin{itemize}
        \item Los alemanes cambiaron periódicamente los patrones de cableado interno de los rotores para 
        evitar que los criptoanalistas aliados pudieran descifrar los mensajes utilizando configuraciones 
        previamente descubiertas.
        \item Además, implementaron cambios regulares en los patrones de las levas utilizadas en el mecanismo 
        de cifrado de Tunny, una máquina de cifrado de teletipo más avanzada. Estos cambios obligaban a 
        los criptoanalistas aliados a adaptar constantemente sus técnicas de descifrado.
    \end{itemize}

    \item \textbf{Desarrollo de Máquinas de Cifrado de Teletipo:}
    \begin{itemize}
        \item La máquina Enigma era una máquina de cifrado mecánica que requería la intervención manual 
        para introducir el texto y transmitir el mensaje cifrado.
        \item Los alemanes desarrollaron máquinas de cifrado de teletipo, como la Lorenz Tunny, que 
        automatizaban el proceso de cifrado y transmisión. Estas máquinas eran más rápidas y eficientes 
        que la Enigma y permitían enviar mensajes más largos. La complejidad de estas máquinas obligó a 
        los aliados a desarrollar nuevas tecnologías de descifrado, como la máquina Colossus.
    \end{itemize}

    \item \textbf{Influencia de las Necesidades de la Guerra:}
    \begin{itemize}
        \item La evolución de la máquina Enigma fue impulsada en gran medida por la presión constante de 
        los criptoanalistas aliados. A medida que los aliados desarrollaban nuevas técnicas para descifrar 
        el código, los alemanes se veían obligados a mejorar la seguridad de sus máquinas.
        \item Las demandas de la guerra, como la necesidad de enviar mensajes más largos y rápidamente, 
        también influyeron en la evolución de la Enigma, llevando al desarrollo de máquinas más sofisticadas.
    \end{itemize}
\end{enumerate}



\section*{Desciframiento de Enigma }

\subsubsection{Antes de Turing}
Durante la década de 1930, Polonia desempeñó 
un papel fundamental en el desciframiento de la 
máquina Enigma, utilizada por el ejército alemán 
durante la Segunda Guerra Mundial. Este esfuerzo 
fue liderado por el Biuro Szyfrów, la oficina de 
cifrado polaca, que logró lo que muchos consideraban 
imposible: romper el código Enigma. Este hito se alcanzó 
gracias a una combinación de brillantez matemática, ingenio 
técnico y la valiosa información proporcionada por agentes 
franceses.

Uno de los protagonistas más destacados en esta hazaña fue 
Marian Rejewski, un matemático polaco cuyo trabajo se volvió 
esencial para el éxito del criptoanálisis polaco. Su equipo 
logró obtener manuales de operación de Enigma que revelaron 
detalles cruciales sobre su funcionamiento interno. Con esta 
información en mano, los criptoanalistas polacos pudieron 
reconstruir una réplica de la máquina y comenzar a analizar 
sus vulnerabilidades. Rejewski se centró en los indicadores, 
grupos de letras transmitidos al inicio de cada mensaje que 
indicaban la configuración de la máquina. Al identificar una 
debilidad en el sistema de indicadores —la repetición diaria 
de configuraciones— logró simplificar el proceso de descifrado 
al darse cuenta de que los primeros caracteres de los 
mensajes cifrados con la misma configuración constituían 
un cifrado de sustitución simple.


Para automatizar y acelerar el proceso de descifrado, los 
polacos desarrollaron las primeras "bombas", máquinas 
electromecánicas capaces de probar miles de configuraciones 
de Enigma en un tiempo relativamente corto. Este avance 
representó un progreso significativo en el campo del 
criptoanálisis y permitió a Polonia interceptar y leer 
mensajes alemanes con regularidad. Gracias a estas 
innovaciones, los criptoanalistas polacos no solo lograron 
desentrañar mensajes cruciales, sino que también sentaron 
las bases para futuras investigaciones en criptografía.

Con la amenaza inminente de la guerra en 1939, Polonia 
decidió compartir sus conocimientos con los aliados 
británicos y franceses. En julio de ese año, se celebró 
una reunión secreta en Varsovia donde los polacos revelaron 
sus técnicas de descifrado, proporcionaron réplicas de Enigma 
y entregaron hojas perforadas utilizadas para el 
criptoanálisis. Esta colaboración fue crucial para 
los esfuerzos británicos en Bletchley Park, donde 
se continuaría el trabajo iniciado por los polacos.A 
pesar del escepticismo inicial por parte de algunos 
aliados sobre la complejidad del código Enigma, la 
información proporcionada por Polonia resultó ser un 
recurso invaluable.

Sin embargo, a pesar del esfuerzo conjunto y las 
contribuciones significativas, los alemanes no se 
quedaron quietos. A lo largo del conflicto, introdujeron 
mejoras en Enigma que complicaron su desciframiento y 
aumentaron su seguridad. Esto presentó un desafío 
adicional para los criptoanalistas aliados, quienes 
debían adaptarse constantemente a las nuevas 
configuraciones y métodos utilizados por el ejército alemán.

A pesar de estas adversidades, el legado del trabajo 
polaco fue invaluable. La investigación realizada por 
Rejewski y su equipo demostró que Enigma era vulnerable y 
que sus mensajes podían ser descifrados. Las técnicas y 
conocimientos compartidos por los criptoanalistas polacos 
inspiraron a figuras como Alan Turing, quien lideraría el 
esfuerzo británico para descifrar los códigos alemanes en 
Bletchley Park. Sin la colaboración polaca, es probable que 
el desciframiento de Enigma se hubiera retrasado 
significativamente, lo que podría haber tenido consecuencias 
desastrosas para el curso de la guerra.

En resumen, Polonia no solo sentó las bases del 
criptoanálisis moderno durante este periodo crítico, 
sino que también contribuyó decisivamente a los esfuerzos 
aliados en la Segunda Guerra Mundial. Su legado perdura no 
solo en la historia del espionaje y la criptografía, sino 
también como un testimonio del ingenio humano frente a 
desafíos aparentemente insuperables. La colaboración entre 
Polonia y sus aliados fue un ejemplo brillante de cómo el 
trabajo conjunto puede cambiar el rumbo de la historia.


\printbibliography


\end{document}

